% Vista editorial exacta de corpus/source/masas/04_persistencia.tex, líneas 365--590. Las líneas 1--364 ya residen exactamente en 73_rev2_electron_lectores_torres.tex.
\subsection{Persistencia del estado enriquecido bajo la extensión superviviente y conservación de la memoria de ruta}
Sea \(x\) un estado APP, \(\tau\) su palabra trítica y \(B_K\) el bloque de memoria al nivel \(K\). El estado relevante para masa tiene la forma

\[
X_K=(x_K,\tau_K,\varepsilon_K,h_K,v_K,c_K,B_K,g_K),
\]
donde \(\varepsilon_K\) es la orientación de hoja, \((h_K,v_K)\) la orientación celular, \(c_K\) el acarreo y \(g_K\) la fase nonádica. Una extensión es superviviente cuando forma una sección compatible de los cilindros sucesivos. Si \(C_{K,j}\), \(0\le j<729\), son los hijos del cilindro en el nivel siguiente, la regla de poda satisface

\[
\sum_{j=0}^{728}G_9(C_{K,j})=1.
\]
El índice único \(j_K\) produce la transición

\[
P_{K+1}=729P_K+j_K.
\]
Al avanzar nueve niveles se conserva la fase,

\[
g(K+9)=g(K),
\]
pero cambian el bloque, el prefijo y la memoria. Por eso el retorno de fase no es retorno de estado. Una partícula estable se define por la compatibilidad de la sección a través de estos retornos, no por la permanencia inmóvil de una coordenada visible.

La ruta observable \(\gamma_{\rm obs}\) es la sombra cronológica de la extensión. El punto que recorre un diagrama es un cursor de lectura; la partícula es la clase completa de persistencia. Esta distinción permite interpretar correctamente el movimiento cuántico: no se exige una bolita siguiendo una trayectoria clásica, sino una historia coherente que admite varias proyecciones compatibles.

\Needspace{7\baselineskip}
\subsection{Registro dodecafásico y firma}
La rueda temporal contiene doce acontecimientos orientados. Su descomposición causal es

\[
12=1+5+6:
\]
un acontecimiento fija el ancla, cinco determinan la incidencia visible y seis forman los pares opuestos que sobreviven al cociente orientado. Si \(s_C^{\rm tot}\) es la suma contraangular entera sobre los doce acontecimientos, el registro conserva

\[
s_C:=s_C^{\rm tot}\in\mathbb Z,
\qquad
s_C^{(12)}:=\frac{s_C^{\rm tot}}6.
\]
El primer objeto es la coordenada entera de la firma; el segundo es su lectura angular dodecafásica. El divisor no es una normalización introducida desde una unidad externa. Es el índice que aparece al reconstruir los seis pares opuestos. La matriz de incidencia dodecafásica posee forma normal de Smith con factores \((1,12)\); el carácter queda bien definido sobre el cociente porque el residuo de orientación se mide en esa estructura.

Ésta es la convención rectora de toda la obra. Las expresiones históricas conservadas que denominaron \(s_C\) a la coordenada ya dividida se leen como \(s_C^{(12)}\); sus cifras y testimonios permanecen intactos, pero ninguna fórmula vigente divide dos veces la misma coordenada.

La pareja de devanados directos es

\[
W_+=6n_A+s_C^{\rm tot},
\qquad
W_-=6n_A-s_C^{\rm tot}.
\]
Como \(s_C=s_C^{\rm tot}\) en la firma publicada, esta fórmula coincide con \(W_\pm=6n_A\pm s_C\). El exponente contraangular usa una sola vez \(s_C^{(12)}=s_C/6\); no vuelve a dividir una coordenada que ya hubiese sido normalizada.

El registro completo publica la firma

\[
\sigma(\Gamma)=
(n_A,s_C,k_{\Delta_4},b_{120},b_\zeta)\in\mathbb Z^5,
\qquad
\zeta_{270}=135\Delta_4^2.
\]
Los cinco componentes tienen tipos distintos. \(n_A\) cuenta el desarrollo angular directo; \(s_C\) registra la suma contraangular entera; \(s_C^{(12)}=s_C/6\) es su lectura angular; \(k_{\Delta_4}\) lee el refinamiento global; \(b_{120}\) marca el canal armónico de altura 120; \(b_\zeta\) pertenece al corrector cuadrático de altura 270. En particular, \(s_{270}\) y \(\zeta_{270}\) no son el mismo objeto: el primero es una traza impar del operador bicapa; el segundo es una coordenada cuadrática del registro.

La palabra electrónica

\[
\tau_{12}=+++---+++---
\]
tiene suma lineal nula, pero su firma par no es nula:

\[
(Q_3,Q_4,Q_6,\rho_0,w)=(-12,-4,12,-6,12).
\]
Tampoco debe identificarse con la firma cruda del ciclo electrónico \((24,0,12,0,-12)\), ni con el origen afín \(v_e=0\). Son cuatro niveles de lectura diferentes: palabra, firma par, firma cruda y origen del carácter relativo.

\Needspace{7\baselineskip}
\subsection{Ángulo, contraángulo y carácter}
Para que la construcción sea autónoma, fijamos aquí las tres coordenadas internas que intervendrán en toda la ley:

\[
\alpha_{\rm HMT}
=\frac{7297352569283800997285105472380663}{10^{36}},
\]
\[
\Delta_4
=\frac{\pi_{\rm HMT}-e_{\rm HMT}\log\pi_{\rm HMT}}{270}
\left(1+\frac{\pi_{\rm HMT}}{729}\right),
\]
y, para \(a>0\) y \(\phi>0\),

\[
H_5(a,\phi)
=\frac12\sqrt{\frac{1000a}{\phi}}-a
+\frac{9a^2}{16}-\frac{5a^3}{9}
+\frac{7a^4}{48}-\frac{a^5}{54}.
\]
En lo sucesivo,

\[
H_5:=H_5(\alpha_{\rm HMT},\varphi_{\rm HMT}).
\]
Estas definiciones no importan una constante metrológica externa: \(\pi_{\rm HMT},e_{\rm HMT},\varphi_{\rm HMT}\) son las evaluaciones coinductivas del lector HMT y \(\alpha_{\rm HMT}\) es la salida racional del cierre dodecafásico. La constante interna \(\alpha_{\rm HMT}\) precede al contraángulo. La rama regional determina primero

\[
A=1000\alpha_{\rm HMT},
\qquad
D_A=\exp\!\left(-\frac{100\pi\alpha_{\rm HMT}}9\right),
\]
\[
\mathcal C_\pi(\alpha_{\rm HMT})
=169\alpha_{\rm HMT}^{6}
+\frac{D_A\alpha_{\rm HMT}^{7}}
       {1-D_A\alpha_{\rm HMT}},
\qquad
y_*=\frac{\pi}{90}(H_5+\alpha_{\rm HMT})
-\mathcal C_\pi(\alpha_{\rm HMT}),
\]
\[
C^*=\frac{180}{\pi}y_*.
\]
Aquí \(A\) y \(C^*\) son coordenadas angulares ya construidas, no magnitudes de acción. Si

\[
\alpha_\angle:=\alpha_{\rm HMT}\,{}^\circ,
\qquad
\eta_{\rm ret}^{(\deg)}:=\frac{C^*}{2}-\alpha_\angle,
\]
entonces

\[
C^*=2\bigl(\eta_{\rm ret}^{(\deg)}+\alpha_\angle\bigr)
\]
es la identidad inversa de retorno, no una segunda regla de selección del contraángulo. Esta escritura evita confundir la mantisa angular \(\eta_{\rm ret}^{(\deg)}\) con la magnitud dimensional \(\hbar_{\rm ret}\). Las cartas en radianes se obtienen mediante la multiplicación por \(\pi/180\). Los canales conjugados son

\[
q_\pm=\exp\left[-\frac\pi{180}(A\pm C^*)\right].
\]
El carácter central de una firma es

\[
\chi_0(\sigma)=
\exp\left(
n_AA_{\rm rad}
+\frac{s_C}{6}C^*_{\rm rad}
+k_{\Delta_4}\Delta_4
+b_\zeta\zeta_{270}
\right).
\]
La coordenada \(b_{120}\) permanece en la firma, pero su contribución se incorpora en la familia de torre. Esta convención impide contar dos veces la misma altura.

El carácter es adimensional. Su función no consiste en convertirse en una constante de acción, sino en actuar por homotecias positivas sobre una recta de masa ya tipada. Para un par de estados \(X,Y\),

\[
\frac{m_Y^{(0)}}{m_X^{(0)}}
=\chi_{\rm full}(\sigma_Y-\sigma_X,\eta_Y-\eta_X).
\]
Esta fórmula separa las razones basales de masa de la elección común de carta absoluta. El cociente beta completo incorpora después \(R_{\rm act}^{\kappa_Y-\kappa_X}\).

\Needspace{7\baselineskip}
\subsection{Semigrupo global de torres de masa y expansión armónica convergente}
Sea \(R^2=I\), con proyectores

\[
P_\pm=\frac{I\pm R}{2},
\qquad
x=\frac{\pi A}{180},
\qquad
y=\frac{\pi C^*}{180},
\]
y

\[
T=e^{-xI-yR}=q_+P_++q_-P_-.
\]
Las trazas par e impar son

\[
c_n=\operatorname{Tr}(T^n)=q_-^n+q_+^n
=2e^{-nx}\cosh(ny),
\]
\[
s_n=-\operatorname{Tr}(RT^n)=q_-^n-q_+^n
=2e^{-nx}\sinh(ny).
\]
Satisfacen

\[
c_n^2-s_n^2=4e^{-2nx},
\]
\[
c_{m+n}=\frac{c_mc_n+s_ms_n}{2},
\qquad
s_{m+n}=\frac{s_mc_n+c_ms_n}{2},
\]
y la recurrencia lineal de orden dos

\[
X_{n+2}=c_1X_{n+1}-e^{-2x}X_n,
\qquad X\in\{c,s\},
\qquad c_0=2,\quad s_0=0.
\]
El conjunto de alturas no se reduce a ocho correcciones visibles. Es el semigrupo

\[
\langle90,120\rangle
=\{90,120\}\cup\{30r:r\ge6\}.
\]
Las alturas \(90,120,180,210,240,270,360,420\) son los primeros testigos genealógicos útiles, no el dominio entero. El carácter refinado es

\[
\log\chi_{\rm full}(\sigma,\eta)
=\log\chi_0(\sigma)
+\sum_{h\in\langle90,120\rangle}N_h(\sigma,\eta)s_h,
\]
con

\[
N_h(\sigma,\eta)=\eta_h+b_{120}\delta_{h,120}.
\]
Así \(s_{120}\) aparece una sola vez, y en altura 270 se conservan separados \(N_{270}s_{270}\) y \(b_\zeta\zeta_{270}\). La convergencia absoluta exige

\[
\sum_h |N_hs_h|<\infty.
\]
La completitud del semigrupo garantiza capacidad representacional de razones positivas; no elige por sí sola los coeficientes físicos de una especie. El regla prospectiva de selección, cuando interviene, debe proceder de ruta, registro genealógico, acarreo, memoria y devanado antes de abrir la masa externa.

\Needspace{7\baselineskip}
